\section{The Tree, in the Order You Meet It}
\label{sec:clocks-tree}

A Series 2 clock tree is not complicated once it is drawn as three questions. This section answers
them for this board only, because the full diagram in the reference manual has forty boxes and no
project needs forty.

\subsection{Question 1: Which Oscillator}

\begin{table}[htbp]
    \centering
    \footnotesize
    \begin{tabular}{@{}p{0.20\textwidth}p{0.16\textwidth}p{0.52\textwidth}@{}}
        \toprule
        \textbf{Oscillator} & \textbf{Rate} & \textbf{When it is the right answer} \\
        \midrule
        \codeinline{HFRCO0} & 19 MHz default & the reset state; enough for GPIO, timers and I2C \\
        \codeinline{HFXO0}  & 38.4 MHz crystal & radio, USB, anything needing ppm accuracy \\
        \codeinline{FSRCO}  & 20 MHz fast RC & cheap and fast, but see the IADC below \\
        \codeinline{HFRCOEM23} & 19 MHz & low-energy domain; the IADC and TIMER0 like it \\
        \codeinline{DPLL0}  & up to 78 MHz & the CPU's fast mode, locked to HFRCO or HFXO \\
        \codeinline{LFRCO}/\codeinline{LFXO} & 32.768 kHz & RTC, sleep, wake-up \\
        \bottomrule
    \end{tabular}
    \caption{Six oscillators, and the four that matter for a board with an I2C bus, a buzzer and a
    microphone.}
    \label{tab:clocks-oscillators}
\end{table}

\subsection{Question 2: Which Branch}

Each oscillator feeds branches, and each branch has a selector and a divider:

\begin{minted}{rust}
// The low-energy peripherals (IADC, TIMER0 in this project) are clocked from EM01GRPACLK.
dp.cmu_s.em01grpaclkctrl().modify(|_r, w| w.clksel().hfrcoem23());
\end{minted}

The name matters when reading the manual: \codeinline{EM01GRPACLK} is one branch,
\codeinline{EM23GRPACLK} another, and a peripheral belongs to exactly one. Choosing the right
oscillator on the wrong branch gives a peripheral clocked at zero — which looks identical to a
peripheral whose registers are not being written.

\subsection{Question 3: Which Gate}

Every peripheral has a bus-clock gate in \codeinline{CMU\_CLKEN0} or \codeinline{CMU\_CLKEN1}. On
this part the GPIO bank is bit 0 of \codeinline{CLKEN0} and the MVP's is bit 30 of
\codeinline{CLKEN1}. A register access with the gate closed is a bus fault, precisely reported:

\begin{minted}{text}
CFSR = 0x00008200   # PRECISERR set: a precise data bus fault
BFAR = 0x5003C040   # ...at PORTA_DOUT, whose clock was off
\end{minted}

That pair of registers is the fastest diagnosis in this book. If a peripheral reads as zero, check
the gate; if a write faults, read \codeinline{CFSR} and \codeinline{BFAR} before anything else.

\subsection{The Three Questions, Together}

\begin{minted}{rust}
// Oscillator: force HFRCOEM23 on and wait for it.
dp.cmu_s.clken0().modify(|_r, w| w.hfrcoem23().set_bit());
dp.hfrcoem23_s.ctrl().modify(|_r, w| w.forceen().set_bit());
while !dp.hfrcoem23_s.status().read().rdy().bit_is_set() {}

// Branch: the low-energy domain comes from it.
dp.cmu_s.em01grpaclkctrl().modify(|_r, w| w.clksel().hfrcoem23());

// Gate: the peripheral's own bus clock.
dp.cmu_s.clken0().modify(|_r, w| w.timer0().set_bit());
\end{minted}

Three decisions, three writes, and one wait. The wait is the subject of the next-but-one section: it
is the only one of the four steps that can hang.
